\documentclass[11pt,letterpaper]{article}

\usepackage[margin=1in]{geometry}
\usepackage{amsmath,amssymb}
\usepackage{enumitem}
\usepackage[hidelinks]{hyperref}

\title{Cost-Constrained Sequential Acquisition of Complete Data Batches}
\author{}
\date{}

\begin{document}
\maketitle

\begin{abstract}
We study the sequential acquisition of complete supervised-learning data.
Each acquisition obtains a batch containing both predictor variables and
target values. Data acquisition, model retraining, and model evaluation
consume separate resources. The system must therefore decide whether to
acquire another batch, retrain the predictive model, evaluate the current
model, or stop. Retraining changes both predictive performance and the utility
assigned to future data acquisitions. We formulate this problem as a
constrained Markov decision process and use Q-learning to learn a sequential
policy over the four actions. We consider two treatments of resource
constraints: hard budget constraints enforced through action masking, and
expected constraints handled through a Lagrangian objective.
\end{abstract}

\section{Problem Setting}

Let \(D^{\mathrm{train}}_0\) denote an initially available labeled training
dataset. Let
\[
    \mathcal{B}_0=\{B_1,\ldots,B_M\}
\]
denote a collection of candidate data batches. Each batch contains complete
supervised observations:
\[
    B_j=(X_j,Y_j),
\]
where \(X_j\) contains predictor variables and \(Y_j\) contains the
corresponding target values.

Before acquisition, the complete contents of a batch are unavailable to the
learning system. The system may observe a low-cost preview or descriptor
\(z_j\), such as source information, batch size, acquisition price,
collection period, coarse feature summaries, or an unlabeled preview sample.
Acquiring a batch reveals both \(X_j\) and \(Y_j\).

The system must sequentially decide:
\begin{enumerate}[label=(\roman*),nosep]
    \item whether to acquire another batch;
    \item when to retrain the model;
    \item when to evaluate the model; and
    \item when to stop.
\end{enumerate}

\section{Experimental Dataset}

The first experiment uses the MiniBooNE Particle Identification dataset from
the UCI Machine Learning Repository. The dataset contains 130,065
observations, 50 continuous predictor variables, and a binary target
distinguishing signal events from background events \cite{miniboone}.

The data are partitioned into:
\begin{itemize}[nosep]
    \item an initially available labeled training set;
    \item an acquisition dataset divided into candidate batches;
    \item an evaluation set;
    \item a hidden validation set; and
    \item a held-out test set.
\end{itemize}

The initial training size, candidate-batch size, number of candidate batches,
and split proportions are experimental parameters. The first prototype uses
100 initial training records and candidate batches of 50 records.

\section{Predictive Model}

The predictive model is a histogram-based gradient-boosted tree classifier.
For an observation \(x\), the model estimates
\[
    f_{\theta_t}(x)=P_{\theta_t}(Y=1\mid x),
\]
where \(\theta_t\) denotes the current model parameters.

Histogram-based gradient boosting is used because it captures nonlinear
relationships, produces class-probability estimates, trains efficiently on
medium-sized tabular datasets, and supports repeated retraining during
sequential experiments. Model hyperparameters are left to experimental
selection.

\section{Batch Acquisition Utility}

Each candidate batch \(B_j\) has a visible descriptor or preview \(z_j\). The
current predictive model assigns an acquisition utility
\[
    q_t(B_j)=g(z_j;\theta_t,m_t),
\]
where \(m_t\) contains the current acquisition and evaluation diagnostics.

A concrete first implementation uses an unlabeled preview sample
\(P_j\subset X_j\) and computes mean predictive entropy:
\[
    q_t(B_j)
    =
    \frac{1}{|P_j|}
    \sum_{x\in P_j}
    H\!\left(P_{\theta_t}(Y\mid x)\right).
\]
For binary classification,
\[
    H\!\left(P_{\theta_t}(Y\mid x)\right)
    =
    -p_t(x)\log p_t(x)
    -(1-p_t(x))\log(1-p_t(x)),
\]
where \(p_t(x)=P_{\theta_t}(Y=1\mid x)\).

The candidate batch with the highest feasible utility is selected:
\[
    B_t^\star
    =
    \arg\max_{B_j\in\mathcal{B}_t}q_t(B_j).
\]
Because the utility depends on \(\theta_t\), retraining changes the ranking of
future candidate batches.

\section{State}

At decision step \(t\), the state is
\[
    s_t
    =
    \left(
        D_t^{\mathrm{train}},
        D_t^{\mathrm{pending}},
        \mathcal{B}_t,
        \theta_t,
        m_t,
        b_t
    \right).
\]

For tabular Q-learning, this state is represented using a compact discretized
feature vector containing:
\begin{itemize}[nosep]
    \item current training-set size;
    \item pending-data size;
    \item number of remaining candidate batches;
    \item remaining acquisition, retraining, and evaluation budgets;
    \item most recent evaluation score;
    \item acquisitions since the last retraining;
    \item retrainings since the last evaluation; and
    \item mean and maximum candidate-batch utility.
\end{itemize}
Continuous variables are discretized into finite intervals.

\section{Actions}

The action space is
\[
    \mathcal{A}
    =
    \{\textsc{Acquire},\textsc{Retrain},\textsc{Evaluate},\textsc{Stop}\}.
\]

\subsection{Acquire}

The feasible candidate batch with the highest acquisition utility is selected:
\[
    B_t^\star=\arg\max_{B_j\in\mathcal{B}_t}q_t(B_j).
\]
The complete batch is added to pending data and removed from the candidate
collection:
\[
    D_{t+1}^{\mathrm{pending}}
    =
    D_t^{\mathrm{pending}}\cup B_t^\star,
    \qquad
    \mathcal{B}_{t+1}
    =
    \mathcal{B}_t\setminus\{B_t^\star\}.
\]
The model is unchanged:
\[
    \theta_{t+1}=\theta_t.
\]

\subsection{Retrain}

All pending data are incorporated:
\[
    D_{t+1}^{\mathrm{train}}
    =
    D_t^{\mathrm{train}}\cup D_t^{\mathrm{pending}},
    \qquad
    D_{t+1}^{\mathrm{pending}}=\varnothing.
\]
The predictive model is refitted:
\[
    \theta_{t+1}
    =
    \operatorname{Train}\!\left(D_{t+1}^{\mathrm{train}}\right).
\]
Retraining changes predictive performance and future acquisition utilities,
but it does not automatically reveal model performance.

\subsection{Evaluate}

The current model is evaluated on the designated evaluation set. The resulting
performance statistic is added to \(m_t\) and becomes visible to the decision
policy. Evaluation changes policy information but does not modify the model or
training data.

\subsection{Stop}

The episode terminates and returns the current predictive model. Pending data
that have not been incorporated through Retrain do not affect the terminal
model.

\section{Sequential Decision Formulation}

The problem is formulated as a Markov decision process:
\[
    \mathcal{M}=(\mathcal{S},\mathcal{A},P,r,\gamma),
\]
where \(\mathcal{S}\) is the state space, \(\mathcal{A}\) is the action space,
\(P(s_{t+1}\mid s_t,a_t)\) is the transition process, \(r_t\) is the reward,
and \(\gamma\) is the discount factor. The policy directly chooses one of the
four actions at every decision step:
\[
    a_t=\pi(s_t).
\]

\section{Q-Learning}

Q-learning estimates the long-term value \(Q(s,a)\) of taking action \(a\) in
state \(s\). After observing a transition
\((s_t,a_t,r_t,s_{t+1})\), the corresponding value is updated as
\[
    Q(s_t,a_t)
    \leftarrow
    Q(s_t,a_t)
    +
    \alpha
    \left[
        r_t
        +
        \gamma\max_{a'}Q(s_{t+1},a')
        -
        Q(s_t,a_t)
    \right],
\]
where \(\alpha\) is the learning rate and \(\gamma\) controls the importance
of future utility.

The learned policy is
\[
    \pi(s)=\arg\max_a Q(s,a).
\]
During training, exploration is introduced using an epsilon-greedy policy.

\section{Reward}

The primary reward is terminal predictive utility. For nonterminal actions,
\[
    r_t=0.
\]
When Stop is selected,
\[
    r_T=U(\theta_T),
\]
where \(U(\theta_T)\) is measured on the hidden validation set during policy
development. This sparse reward allows Q-learning to propagate terminal model
utility backward to earlier decisions.

A small step penalty may optionally discourage unnecessarily long episodes,
but acquisition, retraining, and evaluation costs are treated separately
through constraints.

\section{Cost Model}

The cumulative costs are
\[
    C_T^{\mathrm{acq}}
    =
    \sum_{t<T}c_{\mathrm{acq}}(a_t,s_t),
    \qquad
    C_T^{\mathrm{train}}
    =
    \sum_{t<T}c_{\mathrm{train}}(a_t,s_t),
\]
and
\[
    C_T^{\mathrm{eval}}
    =
    \sum_{t<T}c_{\mathrm{eval}}(a_t,s_t).
\]

Acquisition cost may depend on batch size, source, purchase price, or
collection difficulty. Retraining cost may depend on training-set size,
elapsed time, memory use, processor use, or energy consumption. Evaluation
cost may depend on evaluation-set size, delayed outcome collection, expert
review, or computational effort.

\section{Approach I: Hard Budgets with Action Masking}

Let the initial budgets be
\[
    B_{\mathrm{acq}},\qquad B_{\mathrm{train}},\qquad B_{\mathrm{eval}}.
\]
The remaining budgets are included in the state:
\[
    b_t
    =
    \left(
        b_t^{\mathrm{acq}},
        b_t^{\mathrm{train}},
        b_t^{\mathrm{eval}}
    \right).
\]
Each resource-consuming action subtracts its cost from its corresponding
budget.

At each step, define the feasible action set:
\[
    \mathcal{A}_{\mathrm{feasible}}(s_t)
    =
    \{a\in\mathcal{A}:c(a,s_t)\leq b_t\}.
\]
Q-learning selects only among feasible actions:
\[
    a_t
    =
    \arg\max_{a\in\mathcal{A}_{\mathrm{feasible}}(s_t)}
    Q(s_t,a).
\]
Infeasible actions are masked and cannot be selected.

The optimization problem is
\[
    \max_\pi \mathbb{E}_\pi[U(\theta_T)]
\]
subject to
\[
    C_T^{\mathrm{acq}}\leq B_{\mathrm{acq}},
    \qquad
    C_T^{\mathrm{train}}\leq B_{\mathrm{train}},
    \qquad
    C_T^{\mathrm{eval}}\leq B_{\mathrm{eval}}.
\]
This approach guarantees exact budget compliance. The first prototype uses
unit action costs and hard-budget action masking.

\section{Approach II: Lagrangian Constrained Q-Learning}

The second approach handles expected resource constraints using nonnegative
Lagrange multipliers
\(\lambda_{\mathrm{acq}},\lambda_{\mathrm{train}},\lambda_{\mathrm{eval}}\).
The constrained objective is
\[
    \max_\pi\mathbb{E}_\pi[U(\theta_T)]
\]
subject to expected cost constraints. The modified reward is
\[
    r_t^\lambda
    =
    r_t
    -
    \lambda_{\mathrm{acq}}c_t^{\mathrm{acq}}
    -
    \lambda_{\mathrm{train}}c_t^{\mathrm{train}}
    -
    \lambda_{\mathrm{eval}}c_t^{\mathrm{eval}}.
\]
Q-learning then uses \(r_t^\lambda\).

The multipliers are updated according to observed budget violations. For
example,
\[
    \lambda_{\mathrm{acq}}
    \leftarrow
    \left[
        \lambda_{\mathrm{acq}}
        +
        \eta_\lambda
        \left(C_T^{\mathrm{acq}}-B_{\mathrm{acq}}\right)
    \right]_+,
\]
with analogous updates for retraining and evaluation, where
\([x]_+=\max(x,0)\).

This approach learns a predictive-utility--cost trade-off, but it does not
guarantee that every individual episode satisfies the budgets.

\section{Comparison of Budget Approaches}

Hard-budget action masking should be used when no episode may exceed the
resource limits, costs are known before execution, and compliance must be
guaranteed.

Lagrangian constrained Q-learning should be used when constraints apply on
average, costs may be uncertain, occasional violations are acceptable, and
the policy should learn the utility--cost trade-off.

The first experiment uses hard-budget action masking as the primary
formulation. Lagrangian constrained Q-learning is reserved for a second
experiment.

\section{Experimental Comparison}

The sequential Q-learning policy should be compared against:
\begin{itemize}[nosep]
    \item retraining after every acquisition;
    \item retraining after a fixed number of acquisitions;
    \item evaluation after every retraining;
    \item fixed evaluation intervals;
    \item one final retraining;
    \item random feasible action selection;
    \item hard-budget Q-learning; and
    \item Lagrangian constrained Q-learning.
\end{itemize}

Reported metrics should include final test accuracy, final test AUROC, all
three cumulative costs, numbers of acquisitions, retrainings, and evaluations,
terminal predictive utility, and budget-violation frequency.

\section{Central Research Question}

\begin{quote}
How should a learning system sequentially choose between acquiring a complete
data batch, retraining the model, evaluating the model, and stopping, in order
to maximize terminal predictive utility under acquisition, retraining, and
evaluation constraints?
\end{quote}

\begin{thebibliography}{1}
\bibitem{miniboone}
UCI Machine Learning Repository.
\newblock MiniBooNE Particle Identification Dataset.
\newblock
\url{https://archive.ics.uci.edu/dataset/199/miniboone+particle+identification}.
\end{thebibliography}

\end{document}
